\subsection{Scalar extension and arithmetic transfer}
\label{sec:realization}

\begin{proposition}[Compatibility with extension of scalars]
\label{prop:base-change}
Let \(K\hookrightarrow K'\) be an extension of characteristic-zero fields.
For a squarefree \(P\in K[T]\) and an admissible translation parameter
\(a\in K\), the map \(\widetilde F_G\) of
Theorem~\ref{thm:polynomial-realization}, after coefficientwise extension to
\(K'\), is exactly the map obtained from the same \(P\) and \(a\) over
\(K'\).  Its distinguished fiber satisfies
\[
 \widetilde F_{G,K'}^{-1}(y_{P,a})
 \simeq
 \Spec\bigl((K[T]/(P))\otimes_KK'\bigr)
 \simeq
 \Spec K'[T]/(P).
\]
\end{proposition}

\begin{proof}
Every formula in \eqref{eq:tq}--\eqref{eq:gauge-map}, the output
normalization \(L\), and the target \eqref{eq:main-target} is obtained from
the coefficients of \(P\) and \(a\) by field operations.  The construction
therefore commutes with scalar extension.  The fiber identity is the usual
base-change isomorphism for a polynomial quotient, combined with
Theorem~\ref{thm:polynomial-realization}.
\end{proof}

\begin{remark}[Supplied data versus automatic choices]\label{rem:base-change-choice}
Proposition~\ref{prop:base-change} fixes the presentation \(P\) and
translation \(a\).  The passage from an abstract finite \'etale algebra to a
primitive element, its squarefree minimal polynomial, and an admissible
translation is noncanonical.  Lean implements these existence choices using
classical choice.  Thus the supplied \((P,a)\)-construction commutes with
scalar extension, but no functorial automatically chosen realization is
asserted.
\end{remark}

\begin{corollary}[Number fields as connected full Keller fibers]\label{cor:fields}
Every finite separable field extension \(L/K\) of degree at least three
occurs as a connected full Keller fiber of a Jacobian-one polynomial map
\(\A^3_K\to\A^3_K\) of coordinate degree at most \(6[L:K]+2\).
Given a primitive element and its minimal polynomial, the map is explicit.
\end{corollary}

\begin{proof}
Apply Theorem~\ref{thm:polynomial-realization} to the separable minimal
polynomial of a chosen primitive element.
\end{proof}

\begin{corollary}[Arithmetic data are preserved]\label{cor:transfer}
Let \(A\) be a finite \'etale algebra of rank at least three over a number
field.  Any realization supplied by
Corollary~\ref{cor:finite-etale-realization} has fiber \(\Spec A\) itself.
It therefore preserves:
\begin{itemize}
\item connectedness and field decomposition;
\item real signature;
\item splitting field and Galois action on geometric points;
\item the local algebras \(A\otimes_KK_v\) at every place;
\item for any fixed integral model, factorization types away from its bad primes;
\item the existence or nonexistence of rational and local points.
\end{itemize}
Thus an arithmetic example built as a finite \'etale algebra of rank at least
three can be transferred without alteration to a full Keller fiber.
Proposition~\ref{prop:base-change} gives the corresponding compatibility with
local fields and other scalar extensions.
\end{corollary}
